\subsection{A larger region protected from exterior centres}
\label{sec:protected-shell}

Fix a centre at the origin, a radius $D\ge2$, and let $Y$ be the
other centres of the packing with norm strictly less than $D$.
For $y\in Y$, put
\[
 a_D(y)=\frac{D^2}{D^2+|y|^2-4},\qquad
 K_D^+(Y)=\operatorname{conv}\bigl(B(D/2)\cup
                       \{a_D(y)y:y\in Y\}\bigr).
\]
Here $B(r)$ is the closed ball about the origin.

\begin{prop}[Protection against every exterior centre]\label{prop:protected-shell}
The true Voronoi cell $V_c$ and the cell $V(Y)$ cut out by $Y$
agree on $K_D^+(Y)$. This hull contains both $B(D/2)$ and the
inversion hull $K(Y)$ of \cref{prop:inversion-hull}.
\end{prop}
\begin{proof}
If $|z|\ge D$ is another packing centre, then
$2\langle y,z\rangle\le |z|^2+|y|^2-4$. With $u=|y|^2-4\ge0$,
monotonicity of $t/(t+u)$ gives
\[
 2\langle a_D(y)y,z\rangle
 \le\frac{D^2(|z|^2+u)}{D^2+u}\le |z|^2.
\]
The ball $B(D/2)$ lies in the same half-space by Cauchy--Schwarz.
Convexity yields $K_D^+(Y)\subset H_z$ for each exterior $z$.
Finally $a_D(y)\ge4/|y|^2$, since $(D^2-4)(|y|^2-4)\ge0$.
The inverted point $4y/|y|^2$ is on the segment from $0$ to
$a_D(y)y$. This proves the containment.
\end{proof}

\begin{prop}[A tilted twenty-four-centre family]\label{prop:protected-tilt}
Take $D=\sqrt6$ and suppose $Y$ consists of twenty-four centres
$y_i=2(1+\eta_i)w_i$ with $|w_i|=1$. After rotation and labelling,
let $u_i$ be the unit $D_4$ root directions. Put
$\eta=\max_i\eta_i$ and $\epsilon=\max_i|w_i-u_i|$.
If $0\le\eta\le1/10$, $\eta_i\ge\eta/2$ for every $i$, and
$\epsilon\le\eta/10$, then
\[
 \vol(V(Y)\cap K(Y))\ge 8+6\eta,
 \qquad \vol(V_c)\ge8+9\eta.
\]
Consequently this family satisfies (G), including its equality condition.
\end{prop}
\begin{proof}
Let $Q$ be the regular reference Voronoi cell and
$K_0=\operatorname{conv}\{2u_i\}$. The coordinates in
\cref{sec:root-system} give $\vol Q=8$, $Q\subset B(\sqrt2)$
and $B(\sqrt2)\subset K_0$. For
\[
 s=\frac{1+\min_i\eta_i}{1+\sqrt2\epsilon}
\]
one has $sQ\subset V(Y)$: for $x\in Q$,
$\langle x,w_i\rangle\le\langle x,u_i\rangle+\sqrt2\epsilon
\le1+\sqrt2\epsilon$.

The vertices defining the two hulls in the proposition lie along
$w_i$ at lengths $2/(1+\eta_i)$ and $2k(\eta_i)$, respectively,
where $k(t)=(1+t)/(1+4t/3+2t^2/3)$. Both factors decrease for
$t\ge0$. Both hulls contain the origin. They therefore contain
the hulls of $2b w_i$, with $b=(1+\eta)^{-1}$ and $b=k(\eta)$,
respectively. For any test vector $v$, the support function of
the hull of $2b w_i$ is at least
\[
 b h_{K_0}(v)-2b\epsilon |v|
 \ge b(1-\sqrt2\epsilon)h_{K_0}(v).
\]
Thus the corresponding hull contains $k_*K_0$, with
$k_*=b(1-\sqrt2\epsilon)$.

The part of $sQ$ beyond any one facet of $k_*K_0$ has volume at
most $8(s-k_*)_+^4$. To see this, choose the facet $x_1=\sqrt2 k_*$.
The inequalities $x_1\pm x_j\le\sqrt2 s$ for $2\le j\le4$
bound the section there by a cube of side $2(\sqrt2 s-x_1)$.
Integrating its volume from $\sqrt2 k_*$ to $\sqrt2 s$ gives the
claimed bound, with zero loss when $s\le k_*$. All facets are
equivalent under orthogonal symmetries preserving $Q$ and $K_0$.
Indeed the facets have normals $\pm e_i$ or
$(\pm1,\pm1,\pm1,\pm1)/2$. Signed permutations and an orthogonal
Hadamard matrix divided by $2$ preserve the root system and send
any of these normals to $e_1$. Taking the union over the twenty-four
facets bounds the loss by $192(s-k_*)_+^4$.

Using $\sqrt2<3/2$, one obtains $1+\eta/3\le s\le1+\eta$,
since $(1+\eta/3)(1+3\eta/20)\le1+\eta/2$ in the stated interval.
Also $(1+\eta)^{-1}\ge1-\eta$ and $k(\eta)\ge1-2\eta/5$.
For the latter, multiply by the positive denominator and subtract:
\[
 (1+t)-(1-2t/5)(1+4t/3+2t^2/3)
 =t/15-2t^2/15+4t^3/15\ge0.
\]
It follows that $s-k_*\le43\eta/20$ for $K(Y)$ and
$s-k_*\le31\eta/20$ for $K_D^+(Y)$. Finally
$8(s^4-1)\ge32\eta/3$ and $\eta^4\le\eta/1000$.
Subtracting the respective losses gives the two stated constants,
using $192(43/20)^4/1000<21/5$ and
$192(31/20)^4/1000<6/5$. The protected-hull identity gives the
second inequality for $V_c$. If $\eta=0$, then $\epsilon=0$ and
the configuration is the reference root system.
\end{proof}

This estimate covers a restricted range of simultaneous changes in
direction and radius. It leaves (G) outside that range and the
remaining cases of (C) unresolved.
